% Generated draft: requires mathematical and source review.
\documentclass[11pt]{article}
\usepackage[margin=25mm]{geometry}
\usepackage{fontspec}
\setmainfont{Latin Modern Roman}
\usepackage{amsmath,amssymb,mathtools,unicode-math}
\setmathfont{Latin Modern Math}
\usepackage{enumitem,needspace,xurl,hyperref}
\setlength{\parindent}{0pt}
\setlength{\parskip}{5pt}
\setlength{\emergencystretch}{3em}
\newcommand{\N}{\mathbb{N}}
\newcommand{\Z}{\mathbb{Z}}
\newcommand{\Q}{\mathbb{Q}}
\newcommand{\R}{\mathbb{R}}
\newcommand{\C}{\mathbb{C}}
\newcommand{\F}{\mathbb{F}}
\newcommand{\A}{\mathbb{A}}
\newcommand{\PP}{\mathbb{P}}
\newcommand{\E}{\mathbb{E}}
\newcommand{\eps}{\varepsilon}
\newcommand{\dd}{\,\mathrm{d}}
\begin{document}
\section*{UnsolvedMath ID: 20000435}
\textbf{Source problem code:} catalogue entry
\subsection{Definitions and mathematical statement}
The AIM item concerns smooth complete intersections $X=V(f_1,f_2)\subset\PP^n$ of two quadrics. Its global context takes a number field, with $k=\Q$ for simplicity, and states that for $n\geq8$, $X(\R)\ne\varnothing$ implies $X(\Q)\ne\varnothing$, attributed to Sansuc, Swinnerton-Dyer and Colliot-Th\'el\`ene. It records a circle-method approach for $n\geq9$ (ten variables), attributed to Heath-Brown.

For $n=7$, the stated local question takes two quadratic forms in eight variables over a $p$-adic field $K$ with smooth common zero locus, hence separable determinant polynomial $\det(\lambda f_1+\mu f_2)$. Must some $[\lambda:\mu]\in\PP^1(K)$ yield a form containing three hyperbolic planes? A hyperbolic plane is a split nonsingular binary quadratic form, and containing three means Witt index at least three. Smoothness alone is the printed hypothesis; the supplied report emphasizes that the intended known local theorem requires $X(K)\ne\varnothing$ as an additional essential assumption.

The report attributes to Heath-Brown the global Hasse principle and weak approximation for smooth intersections in $\PP^7$, and to Colliot-Th\'el\`ene's Th\'eor\`eme 7.3 the intended local statement for every $p$-adic field, including dyadic fields, under that nonemptiness hypothesis. Hasse principle means that points over all completions imply a rational point; weak approximation means density in the product over any finite set of completions. Full field and theorem hypotheses are not supplied, and these attributions do not prove the literal smoothness-only local version.

Its additional local proposition uses an affine pencil $q(t)$ with a $K$-rational simple determinant root $t_0$. Let $\pi$ be a uniformizer. It claims a threshold $M$ and parity $\epsilon\in\{0,1\}$ such that
\[
 \text{for all }m\geq M\text{ with }m\equiv\epsilon\pmod2,
 \quad\text{for all }u\in\mathcal O_K^\times,
 \quad i_W(q(t_0+u\pi^m))=3,
\]
and every indicated form is nonsingular. This is exact index three, rather than merely at least three. The root has rank seven in the separable-pencil setting. A simple root gives determinant valuation $v_K(\det'(q)(t_0))+m$ sufficiently near it; the report does not supply its explicit threshold or parity formula, so the proposition is retained without guessing those constants. It covers this rational-root sector, not every pencil.

The original remarks additionally report an unwritten odd-characteristic odd-degree-extension argument by Heath-Brown, a rank-at-most-seven pencil-member case, and lower-dimensional conditions such as two conjugate lines or a ground-field conic for $n\geq5$. Their exact residue-characteristic, point and dimension hypotheses are incomplete. They remain contextual claims requiring verification.
\subsection{Short English statement}
Ask for a three-hyperbolic-plane member in an eight-variable $p$-adic quadratic pencil. The report identifies the essential local-solubility hypothesis in the known theorem, including dyadic fields, and cites the global Hasse-principle result. Its additional statement gives whole parameter annuli of exact Witt index three near a rational simple determinant root, with one valuation parity.
\subsection{Sources}
\begin{itemize}
\item Local registry: \url{data/unsolvedmath/problems.json}, permanent ID $20000435$, source code \texttt{catalogue entry}; original question, all remarks and supplied research summary.
\item AIM workshop ``Rational and integral points on higher dimensional varieties'', Question 3: \url{https://aimath.org/WWN/qptsurface2/qptsurface2.pdf}. Canonical location \texttt{aim-algebraic-number-theory-notes.json}, \texttt{notes[86]}; original literature field empty. Supplied verification reports a PDF check, not independently performed here.
\item Source attributions: Sansuc, Swinnerton-Dyer, Colliot-Th\'el\`ene and Heath-Brown. Summary additionally cites Colliot-Th\'el\`ene, Th\'eor\`eme 7.3, and Heath-Brown's global result without full bibliographic details.
\item Supplied machine-generated research attempt 1: partial result; problem partially solved at the run and in the registry; rigor metadata ``proved'', novelty confidence low. The non-peer-reviewed annulus proposition requires independent verification before citation.
\end{itemize}
\end{document}
